\documentclass[12pt, twoside]{article}
\input{../preamble}

\fancyhead[LO]{La Scuola d'Italia / Dr.\ Huson / IB Math 12 \\* 14 September 2026}
\fancyhead[RO]{\fbox{\begin{minipage}{6cm}\footnotesize Nickname:\\[0.5cm]\end{minipage}}}
\fancyfoot[C]{\thepage}

\begin{document}
\subsubsection*{1.2 Classwork: Sampling and Histograms}

\begin{enumerate}[itemsep=0.15cm]

%--- Problem 1: Sampling terminology and survey design [7 marks] ---
%--- Source: Original (freshly authored, AA SL 4.1); style after Haese Core Topics SL, Ch 11A ---
\item[1.]

There are 420 students at a school. The principal wants to know what
the students think of the new lunch menu. She stands at the cafeteria
door one Tuesday and asks the first 30 students who walk in:

\begin{center}
\emph{``Do you agree that the new lunch menu is a big improvement?''}
\end{center}

\begin{enumerate}[itemsep=0.1cm]
  \item Write down the population, and write down the size of the sample.
  \item The principal took the students who happened to be closest to
  hand. Name this sampling technique.
  \item Give one reason why her sample may not be representative of all
  420 students.
  \item The question she asked is a \emph{leading} question. Explain what
  is wrong with it, and rewrite it so that it is not leading.
  \item Describe how the principal could instead take a simple random
  sample of 30 students.
\end{enumerate}

\begin{tikzpicture}
  \draw (-0.6,0) rectangle (14.6,12.6);
  \foreach \y in {1,2,...,11} \draw[dotted] (0,\y)--(14.1,\y);
\end{tikzpicture}

\newpage

%--- Problem 2: Grouping continuous data and drawing a histogram [7 marks] ---
%--- Source: Original (freshly authored, AA SL 4.1, 4.3); style after Haese Core Topics SL, Ch 11E ---
\item[2.]

The height of each of 24 tomato plants was measured, in centimetres:

{\small
\begin{center}
\setlength{\tabcolsep}{8pt}
\begin{tabular}{cccccc}
23.4 & 31.7 & 28.0 & 36.2 & 41.5 & 27.3 \\
29.8 & 25.0 & 39.1 & 30.0 & 22.6 & 33.4 \\
37.9 & 29.5 & 26.8 & 43.2 & 35.0 & 31.1 \\
24.7 & 32.6 & 28.9 & 40.3 & 26.1 & 42.0 \\
\end{tabular}
\end{center}}

\begin{enumerate}[itemsep=0.1cm]
  \item Complete the grouped frequency table below.
  \item Write down the modal class.
  \item One plant measured exactly $30.0$ cm. State which class it is
  counted in, and explain why.
  \item Draw a histogram for the data on the grid below.
\end{enumerate}

{\renewcommand{\arraystretch}{1.7}
\begin{center}
\begin{tabular}{|c|p{4.2cm}|c|}
\hline
\textbf{Height $h$ (cm)} & \hfil \textbf{Tally} & \textbf{Frequency} \\ \hline
$20 \le h < 25$ & & \\ \hline
$25 \le h < 30$ & & \\ \hline
$30 \le h < 35$ & & \\ \hline
$35 \le h < 40$ & & \\ \hline
$40 \le h < 45$ & & \\ \hline
\textbf{Total}  & & \\ \hline
\end{tabular}
\end{center}}

\vspace{-0.2cm}

\begin{center}
\begin{tikzpicture}[x=0.44cm, y=0.66cm]
  % grid
  \draw[gray!30] (20,0) grid[xstep=1, ystep=1] (45,9);
  \draw[gray!60, thin] (20,0) grid[xstep=5, ystep=2] (45,9);
  % axes
  \draw[thick, -{Stealth[length=2mm]}] (20,0) -- (46.2,0);
  \node[anchor=west] at (46.6,0) {\footnotesize $h$ (cm)};
  \draw[thick, -{Stealth[length=2mm]}] (20,0) -- (20,9.6);
  \foreach \x in {20,25,30,35,40,45} \draw (\x,0) -- (\x,-0.25) node[below] {\footnotesize \x};
  \foreach \y in {0,2,4,6,8} \draw (20,\y) -- (19.6,\y) node[left] {\footnotesize \y};
  \node[rotate=90, anchor=south] at (17.6,4.5) {\footnotesize frequency};
\end{tikzpicture}
\end{center}

\end{enumerate}
\end{document}
